% Generated draft: requires mathematical and source review.
\documentclass[11pt]{article}
\usepackage[margin=25mm]{geometry}
\usepackage{fontspec}
\setmainfont{Latin Modern Roman}
\usepackage{amsmath,amssymb,mathtools,unicode-math}
\setmathfont{Latin Modern Math}
\usepackage{enumitem,needspace,xurl,hyperref}
\setlength{\parindent}{0pt}
\setlength{\parskip}{5pt}
\setlength{\emergencystretch}{3em}
\newcommand{\N}{\mathbb{N}}
\newcommand{\Z}{\mathbb{Z}}
\newcommand{\Q}{\mathbb{Q}}
\newcommand{\R}{\mathbb{R}}
\newcommand{\C}{\mathbb{C}}
\newcommand{\F}{\mathbb{F}}
\newcommand{\A}{\mathbb{A}}
\newcommand{\PP}{\mathbb{P}}
\newcommand{\E}{\mathbb{E}}
\newcommand{\eps}{\varepsilon}
\newcommand{\dd}{\,\mathrm{d}}
\begin{document}
\section*{UnsolvedMath ID: 20000239}
\textbf{Source problem code:} catalogue entry
\subsection{Definitions and mathematical statement}
Determine the connection between Euler class groups and Chow--Witt groups. In a smooth affine $d$-dimensional setting $X=\operatorname{Spec}A$, the top Euler class group $E^d(A)$ records oriented codimension-$d$ local complete intersections modulo global complete-intersection relations. Chow--Witt groups retain quadratic orientation data; with a line-bundle twist $L$, their usual smooth characteristic-not-two description is $\widetilde{\operatorname{CH}}^r(X,L)=H^r(X,\mathbf K_r^{MW}(L))$. The original does not fix smoothness, affine hypotheses, dimension, field or twists.

Jean Fasel's accompanying question additionally asks for a product-compatible integral real cycle map lifting the mod-two cycle map. The source recalls product-compatible maps $\operatorname{CH}^r(X)\to H^{2r}(X(\C),\Z)$ over $\C$ and $\operatorname{CH}^r(X)\to H^r(X(\R),\Z/2)$ over $\R$. It asks for a map from $\widetilde{\operatorname{CH}}^r(X,\omega_{X/\R})$ to $H^r(X(\R),\Z)$ commuting with forgetful Chow and mod-two reduction maps.

The attached unreviewed aid reports an isomorphism
\[
E^d(A)\simeq\widetilde{\operatorname{CH}}^d(\operatorname{Spec}A)
\]
for smooth affine $d$-folds over infinite fields of characteristic not two. It reports that on a smooth quasi-projective real variety the natural integral realization with twist $L$ instead has target
\[
H^r(X(\R),\Z(L)),
\]
where $\Z(L)$ is the sign local system of the real line bundle $L$. For $L=\omega_{X/\R}$ this is the orientation local system, so a constant-$\Z$ formulation needs an orientation trivialization. Products naturally multiply twists: the $L,M$ groups map to the $L\otimes M$ group. Thus product compatibility should be stated for a family of twists; two $\omega$-twisted inputs do not automatically have an $\omega$-twisted product. The precise realization normalization and comparison theorem hypotheses are not supplied.

As a test family the aid takes $A_d=\R[x_0,\ldots,x_d]/(\sum x_i^2-1)$ and its tangent module $P=\ker(A_d^{d+1}\to A_d)$, where $(v_i)\mapsto\sum x_iv_i$, with determinant orientation. It claims real Euler image $\pm\chi(S^d)$; for even $d\ge2$ this detects an infinite-order element in $\ker(E^d(A_d)\to E_0^d(A_d))$, where $E_0^d$ forgets orientation. For odd $d$ it reports an algebraic unit vector splitting $P$ and vanishing Euler class, but the vector is not given. These are aid claims, not a proof or unrestricted comparison theorem in this draft.
\subsection{Short English statement}
Compare Euler classes with quadratic Chow--Witt classes and construct the real integral cycle map asked for by Jean Fasel. Orientation local systems and multiplication of line-bundle twists are necessary to make the original diagram and product compatibility precise.
\subsection{Sources}
\begin{itemize}
\item AIM \emph{Projective modules and $\A^1$-homotopy theory}, Question 2, ``by popular demand,'' and Jean Fasel's accompanying remark/question: \url{https://aimath.org/pastworkshops/projectiveA1problems.pdf}. Canonical location: \texttt{aim-algebraic-geometry-notes.json}, notes[238]. The record is labeled exact but visibly damaged; its verification references page 2 without reproducing a complete corrected text. The supplied extraction is used without consulting a PDF.
\item The attached aid reports the top-codimension isomorphism, twisted real realization and sphere family without complete bibliographic references. It is machine generated and not peer reviewed, with low novelty confidence; the original literature field is empty.
\item Local registry: \texttt{data/unsolvedmath/problems.json}, permanent catalogue entry, exact source code \texttt{catalogue entry}.
\end{itemize}
\end{document}
